\documentclass[11pt]{article}
\usepackage[margin=1in]{geometry}
\usepackage{amsmath,amssymb}
\usepackage[hidelinks]{hyperref}

\title{Literature Status: Bounded-Gap Quasi-Greediness\\
for Markushevich Bases}
\author{Run \texttt{fa\_banach\_001}, lane 3}
\date{19 July 2026}

\begin{document}
\maketitle

\paragraph{Status.}
\texttt{literature\_already\_answered}.

\paragraph{Original question.}
Berasategui and Bern\'a's \emph{Quasi-greedy bases for sequences with gaps}
\cite{BB21}, arXiv:2005.07221, proves its bounded-quotient-gap implication
under a Schauder-basis assumption.  Section 8, Question 1 (page 17 of the
18-page arXiv PDF) asks:
\begin{quote}
Theorem 5.2 is proven under the condition of Schauder bases. Can the Schauder
condition be removed and the result be obtained using only the Markushevich
condition?
\end{quote}

\paragraph{Identification of the answer.}
The same authors' later paper \emph{Greedy-like bases for sequences with gaps}
\cite{BB24}, arXiv:2009.02257, gives the affirmative answer.  Its Theorem 2.11
(page 8 of the arXiv PDF) states:
\begin{quote}
Let $\mathbf n$ be a sequence with bounded quotient gaps. If $\mathcal B$ is
an $\mathbf n$-quasi-greedy Markushevich basis for a Banach space $X$, it is
quasi-greedy.
\end{quote}
This removes exactly the Schauder hypothesis asked about in the source.  It
also yields the version with a fixed weakness parameter $0<t\leq1$: every
$\mathbf n$-$t$-quasi-greedy system is $\mathbf n$-quasi-greedy, since every
greedy set is a $t$-greedy set.

The supporting paper explicitly says that it generalizes the earlier result
to Markushevich bases and completes the question.  The relation is therefore
an acknowledged later-literature answer, not a new result of this run.

\paragraph{Scope.}
The source's final Question 1 is fully answered.  Other questions involving
arbitrary quotient gaps, weak-parameter converses, or the later paper's new
partial-greediness variants are outside this identification.

\begin{thebibliography}{9}
\bibitem{BB21}
M. Berasategui and P. M. Bern\'a,
\emph{Quasi-greedy bases for sequences with gaps},
Nonlinear Analysis \textbf{208} (2021), 112294;
\href{https://arxiv.org/abs/2005.07221}{arXiv:2005.07221}.

\bibitem{BB24}
M. Berasategui and P. M. Bern\'a,
\emph{Greedy-like bases for sequences with gaps},
Banach Journal of Mathematical Analysis \textbf{18} (2024), article 2;
\href{https://arxiv.org/abs/2009.02257}{arXiv:2009.02257}.
\end{thebibliography}

\end{document}

